%% toc: P2 12 12.5,12.6
\section{Cut, the subformula property, and translations between the calculi}
\label{sec:p02-w02b-p02-009-cut-translations}

This section completes the metatheory of the two classical propositional
calculi of \cref{sec:prop-deduction-02} and of the natural-deduction calculus
\(\mathsf{ND}_{\mathsf{CPL}}\) defined before it.  It has three purposes.
First, it proves that the cut rule of \(\mathsf{LK}_{\mathsf{CPL}}\) is
admissible: every derivable sequent has a derivation without Cut, and in fact
one that uses only initial sequents and logical rules.  Second, it defines
the subformula property and proves that every Cut-free derivation has it.
Third, it gives two explicit translations of derivations, from
\(\mathsf{ND}_{\mathsf{CPL}}\) to \(\mathsf{LK}_{\mathsf{CPL}}\) and back, rule
by rule, with an exact account of which labelled assumption occurrences each
rule discharges and of how weakening and contraction produce the
multiplicities of the sequents.  The rule-by-rule soundness cases of both
calculi, which the soundness theorems \cref{thm:prop-nd-soundness,thm:prop-lk-soundness}
use, are stated here one rule at a time.

\emph{Inputs.}  An \(\mathsf{LK}_{\mathsf{CPL}}\)-derivation, possibly with
Cut; or a valid sequent; or an \(\mathsf{ND}_{\mathsf{CPL}}\)-derivation.
\emph{Conclusions.}  A derivation of the same sequent using only initial
sequents and logical rules; the subformula property of every Cut-free
derivation; an \(\mathsf{LK}_{\mathsf{CPL}}\)-derivation of the sequent whose
antecedent is the multiset of open assumption occurrences of a given
natural-deduction derivation; and a natural-deduction derivation of
\(\bottom\) from the antecedent formulas and the negated succedent formulas
of a given derivable sequent.

\subsection{Setting}

\begin{hbconvention}{def:p02-w02b-p02-009-setting}{The two calculi}
Formulas are the members of \(\mathsf{Fm}_{\mathsf{CPL}}\) defined in
\cref{sec:prop-syntax-semantics-01}, with primitives
\(p_n,\bottom,\neg,\land,\lor,\to\); \(c(\varphi)\) is the complexity and
\(\operatorname{Sub}(\varphi)\) the subformula set defined there, and
valuations and \(\entails\) are those of \cref{sec:prop-syntax-semantics-02}.  Sequents
\(\Gamma\Rightarrow\Delta\) have finite multisets as antecedent and
succedent, the comma is pointwise sum, and \(\operatorname{supp}\) is the
support, as in the definition of sequents in \cref{sec:prop-deduction-02}.
For a finite multiset \(\Pi\) write \(|\Pi|\defeq\sum_{\chi}\Pi(\chi)\) and
\(\varphi^{k}\) for \(k\) copies of \(\varphi\) (the empty multiset when
\(k=0\)).

\(\mathsf{LK}_{\mathsf{CPL}}\) is exactly the calculus displayed in
\cref{sec:prop-deduction-02}: identity sequents
\(\Gamma,\varphi\Rightarrow\varphi,\Delta\) and \(\bottom\)-sequents
\(\Gamma,\bottom\Rightarrow\Delta\) as initial sequents; the structural
rules \(W_L,W_R,C_L,C_R\); the logical rules
\(\neg L,\neg R,\land L,\land R,\lor L,\lor R,\to L,\to R\), whose
two-premise instances share their contexts; and Cut, with premises
\(\Gamma\Rightarrow\Delta,\varphi\) and \(\Sigma,\varphi\Rightarrow\Pi\) and
conclusion \(\Gamma,\Sigma\Rightarrow\Delta,\Pi\).  In a logical rule the
displayed compound formula of the conclusion is \emph{principal} and the
displayed formulas of the premises are its \emph{side formulas}.

\(\mathsf{ND}_{\mathsf{CPL}}\) is exactly the calculus of labelled-discharge
trees displayed before \cref{lem:nd-substitution}, with rules
\(\land I,\land E_1,\land E_2,\lor I_1,\lor I_2,\lor E,\to I,\to E,
\neg I,\neg E,\bottom E,\mathsf{RAA}\).  An application of
\(\to I\), \(\neg I\) or \(\mathsf{RAA}\) with label \(i\) discharges exactly
the open occurrences marked \(i\) of its bracketed formula above its premise,
and \(\lor E\) with labels \(i,j\) discharges exactly the open occurrences of
the first disjunct marked \(i\) above its second premise and of the second
disjunct marked \(j\) above its third premise; every label carried by such an
application discharges at least one occurrence.  A leaf written
\(\varphi^{\ell}\) is an assumption occurrence of \(\varphi\) marked
\(\ell\).
\end{hbconvention}

\begin{hbdefinition}{def:p02-w02b-p02-009-validity}{Falsifying valuations, validity, Cut-free derivations and weight}
A valuation \(v\) \emph{falsifies} a sequent \(\Gamma\Rightarrow\Delta\) iff
\(v\entails\gamma\) for every \(\gamma\in\operatorname{supp}(\Gamma)\) and
\(v\not\entails\delta\) for every \(\delta\in\operatorname{supp}(\Delta)\);
otherwise \(v\) satisfies it, in the sense of the definition of sequents in
\cref{sec:prop-deduction-02}.  The sequent is \emph{valid} iff no valuation
falsifies it.  An \(\mathsf{LK}_{\mathsf{CPL}}\)-derivation is
\emph{Cut-free} iff it contains no application of Cut; it may contain
weakening and contraction.  It is \emph{purely logical} iff each of its nodes
is an initial sequent or the conclusion of a logical rule.  The
\emph{weight} of a sequent is
\[
 \operatorname{wt}(\Gamma\Rightarrow\Delta)\defeq
 \sum_{\chi\in\mathsf{Fm}_{\mathsf{CPL}}}
 \bigl(\Gamma(\chi)+\Delta(\chi)\bigr)\,c(\chi),
\]
a finite sum because \(\Gamma\) and \(\Delta\) have finite support.
\end{hbdefinition}

\subsection{Rule-by-rule soundness}

\begin{hblemma}{thm:p02-w02b-p02-009-lk-rule-soundness}{proved=annex:p02-w02b-p02-009-lk-rule-soundness}{Every rule of the sequent calculus reflects falsification}
No valuation falsifies an initial sequent of
\(\mathsf{LK}_{\mathsf{CPL}}\).  For every instance of each of the rules
\(W_L,W_R,C_L,C_R,\neg L,\neg R,\land L,\land R,\lor L,\lor R,\to L,\to R\)
and Cut, every valuation that falsifies the conclusion falsifies at least one
premise.  Consequently every sequent derivable in
\(\mathsf{LK}_{\mathsf{CPL}}\) is valid; this is
\cref{thm:prop-lk-soundness}, here with every rule case written out.
\end{hblemma}

\begin{hblemma}{thm:p02-w02b-p02-009-nd-rule-soundness}{proved=annex:p02-w02b-p02-009-nd-rule-soundness}{Every rule of natural deduction preserves truth from open assumptions}
Let \(D\) be an \(\mathsf{ND}_{\mathsf{CPL}}\)-derivation with conclusion
\(\varphi\), and let \(v\) be a valuation that satisfies the formula of
every open assumption occurrence of \(D\).  Then \(v\entails\varphi\).  The
proof treats the assumption case and each of the twelve rules separately;
for a set \(\Gamma\) it gives
that \(\Gamma\derives_{\mathsf{ND}_{\mathsf{CPL}}}\varphi\) implies
\(\Gamma\entails\varphi\), which is \cref{thm:prop-nd-soundness}.
\end{hblemma}

\subsection{Cut admissibility and the subformula property}

The route to a Cut-free derivation used here is semantic: each logical rule
is invertible with respect to validity, so a valid sequent can be decomposed
from the root upward until only initial sequents remain, and soundness
places every derivable sequent among the valid ones.  This route is the one
Buss uses for his propositional calculus \(PK\)
\cite[Secs.~1.2.7--1.2.11, pp.~13--16]{Buss1998Intro}; the statements below
are proved directly for \(\mathsf{LK}_{\mathsf{CPL}}\), whose multiset
contexts, contexted identity sequents, \(\bottom\)-sequents and primitive
negation differ from \(PK\) (\cref{def:p02-w02b-p02-009-scope}).

\begin{hblemma}{thm:p02-w02b-p02-009-inversion}{proved=annex:p02-w02b-p02-009-inversion}{Semantic invertibility of the logical rules}
For every instance of each logical rule of \(\mathsf{LK}_{\mathsf{CPL}}\),
every valuation that falsifies a premise falsifies the conclusion.  Hence, if
the conclusion is valid, every premise is valid.  Moreover, the weight of
each premise is strictly smaller than the weight of the conclusion.
\end{hblemma}

\begin{hbtheorem}{thm:p02-w02b-p02-009-cutfree-completeness}{proved=annex:p02-w02b-p02-009-cutfree-completeness}{Purely logical completeness}
Every valid sequent \(\Gamma\Rightarrow\Delta\) has a purely logical
\(\mathsf{LK}_{\mathsf{CPL}}\)-derivation: a derivation using only initial
sequents and the eight logical rules, with no Cut, weakening or
contraction.
\end{hbtheorem}

\begin{hbtheorem}{thm:p02-w02b-p02-009-cut-admissibility}{proved=annex:p02-w02b-p02-009-cut-admissibility}{Cut admissibility for the classical propositional sequent calculus}
If \(\Gamma\Rightarrow\Delta\) is derivable in \(\mathsf{LK}_{\mathsf{CPL}}\),
with any number of applications of Cut, weakening and contraction, then
\(\Gamma\Rightarrow\Delta\) has a purely logical derivation.  In particular,
Cut is admissible: adding it to the other rules of
\(\mathsf{LK}_{\mathsf{CPL}}\) derives no new sequent; and
\(W_L,W_R,C_L,C_R\) and Cut are together admissible in the calculus of
initial sequents and logical rules.  For every finite multisets
\(\Gamma,\Delta\), the following are equivalent: \(\Gamma\Rightarrow\Delta\)
is derivable; it has a purely logical derivation; it is valid.
\end{hbtheorem}

The theorem asserts the existence of a Cut-free derivation of the same
endsequent; it does not transform a given derivation.  The derivation it
provides is the one built by the procedure of
\cref{def:p02-w02b-p02-009-decomposition} from the endsequent alone.

\begin{hbdefinition}{def:p02-w02b-p02-009-subformula-property}{Subformula property}
An \(\mathsf{LK}_{\mathsf{CPL}}\)-derivation with endsequent
\(\Gamma\Rightarrow\Delta\) has the \emph{subformula property} iff for every
sequent \(\Pi\Rightarrow\Sigma\) occurring in it and every
\(\theta\in\operatorname{supp}(\Pi)\cup\operatorname{supp}(\Sigma)\) there is
\(\chi\in\operatorname{supp}(\Gamma)\cup\operatorname{supp}(\Delta)\) with
\(\theta\in\operatorname{Sub}(\chi)\).
\end{hbdefinition}

\begin{hbtheorem}{thm:p02-w02b-p02-009-subformula}{proved=annex:p02-w02b-p02-009-subformula}{Subformula property without Cut}
Every Cut-free \(\mathsf{LK}_{\mathsf{CPL}}\)-derivation, with or without
weakening and contraction, has the subformula property.  More precisely,
for every application of a rule other than Cut, every formula in the
support of either side of a premise is a subformula of a formula in the
support of one side of the conclusion.
\end{hbtheorem}

\begin{hbcorollary}{thm:p02-w02b-p02-009-analytic}{proved=annex:p02-w02b-p02-009-analytic}{Analytic derivations}
Every \(\mathsf{LK}_{\mathsf{CPL}}\)-derivable sequent has a purely logical
derivation with the subformula property.  Every formula occurring in such a
derivation of \(\Gamma\Rightarrow\Delta\) belongs to the finite set
\(\bigcup\{\operatorname{Sub}(\chi):\chi\in\operatorname{supp}(\Gamma)\cup
\operatorname{supp}(\Delta)\}\).  The empty sequent \({}\Rightarrow{}\) is
not derivable.
\end{hbcorollary}

\begin{hbmethod}{def:p02-w02b-p02-009-decomposition}{Root-first decomposition}
\emph{Input:} a sequent \(\Gamma\Rightarrow\Delta\).  \emph{Output:} a purely
logical derivation of it, or a valuation falsifying it.
\begin{enumerate}
  \item If \(\bottom\in\operatorname{supp}(\Gamma)\), stop: the sequent is a
    \(\bottom\)-sequent.  If some \(\varphi\) lies in
    \(\operatorname{supp}(\Gamma)\cap\operatorname{supp}(\Delta)\), stop: it
    is an identity sequent.
  \item Otherwise, if every member of
    \(\operatorname{supp}(\Gamma)\cup\operatorname{supp}(\Delta)\) is a
    variable or \(\bottom\), stop and return the valuation \(v\) with
    \(v(p_n)=1\) iff \(p_n\in\operatorname{supp}(\Gamma)\).
  \item Otherwise choose one occurrence of a compound formula, write the
    sequent as the conclusion of the logical rule whose principal formula is
    that occurrence (unique readability fixes the rule), apply the procedure
    to each premise, and place the rule beneath the results.  If a premise
    returns a falsifying valuation, return it.
\end{enumerate}
\emph{Side conditions:} the contexts of the premises are the remaining
contexts of the conclusion, shared in \(\land R\), \(\lor L\), \(\to L\).
Every call terminates, because premises have smaller weight
(\cref{thm:p02-w02b-p02-009-inversion}).  A returned valuation falsifies the
premise it came from and therefore the input sequent, by the same lemma;
the stop in step 2 returns a falsifying valuation because neither stop of
step 1 applied.
\end{hbmethod}

\begin{hbexample}{def:p02-w02b-p02-009-peirce}{Peirce's law without Cut}
Decomposing \({}\Rightarrow((P\to Q)\to P)\to P\) gives, in three logical
steps, the following derivation, written as numbered lines in which each
line names the rule and the lines of its premises:
\[
\begin{array}{rll}
1.&P\Rightarrow P,Q&\text{identity sequent}\\
2.&{}\Rightarrow P,P\to Q&\to R,\ 1\\
3.&P\Rightarrow P&\text{identity sequent}\\
4.&(P\to Q)\to P\Rightarrow P&\to L,\ 2,3\\
5.&{}\Rightarrow((P\to Q)\to P)\to P&\to R,\ 4
\end{array}
\]
Both leaves are identity sequents, the second with empty contexts and the
first with \(\Delta=Q\).  The \(\to L\) step has \(\Gamma\) empty,
\(\Delta=P\), \(\varphi=P\to Q\) and \(\psi=P\).  Every formula is a
subformula of the endformula.  The left premise of \(\to L\) has two
succedent formulas; the multi-formula succedent is what lets the classical
law have a Cut-free derivation.
\end{hbexample}

\begin{hbwarning}{def:p02-w02b-p02-009-cut-limits}{What Cut changes, and what the semantic proof does not give}
Cut is the only rule whose premises may contain a formula that is not a
subformula of its conclusion.  For instance, Cut on \(R\) between the
identity sequents \(P\Rightarrow P,R\) and \(R,P\Rightarrow P\) concludes
\(P,P\Rightarrow P,P\), and \(C_L\), \(C_R\) then give \(P\Rightarrow P\);
the derivation contains \(R\), which is not a subformula of \(P\).

The proof of \cref{thm:p02-w02b-p02-009-cut-admissibility} discards the
given derivation and rebuilds one from the endsequent.  It yields no
procedure acting on derivations and no bound on the size of the result in
terms of the size of the given derivation; Buss records both points for
the same semantic route \cite[Sec.~1.2.11, pp.~15--16]{Buss1998Intro}.
\end{hbwarning}

\begin{hbconvention}{def:p02-w02b-p02-009-scope}{Exact scope of this result and of the Part~7 elimination theorem}
\Cref{thm:p02-w02b-p02-009-cut-admissibility} concerns exactly
\(\mathsf{LK}_{\mathsf{CPL}}\).  The syntactic elimination theorem
\cref{thm:p07-classical-cut-elimination} concerns exactly the calculus
\texttt{sheet:lk-defneg}, and Buss's semantic argument concerns his \(PK\)
\cite[Sec.~1.2.2, pp.~11--12]{Buss1998Intro}.  They differ as follows.
\begin{center}
\small
\begin{tabular}{@{}p{0.2\linewidth}p{0.24\linewidth}p{0.24\linewidth}p{0.24\linewidth}@{}}
\toprule
 & \(\mathsf{LK}_{\mathsf{CPL}}\) & \texttt{sheet:lk-defneg} & \(PK\) \\
\midrule
Formulas & \(p_n,\bottom,\neg,\land,\lor,\to\) &
 \(\bot,\land,\lor,\to,\forall,\exists\); \(\neg\varphi\defeq\varphi\to\bot\) &
 variables with \(\neg,\land,\lor,\supset\) \\
Contexts & finite multisets & finite multisets & sequences, explicit exchange \\
Identity & \(\Gamma,\varphi\Rightarrow\varphi,\Delta\), any \(\varphi\) &
 \(\varphi\Rightarrow\varphi\) & \(A\Rightarrow A\), \(A\) atomic \\
Falsum & \(\Gamma,\bottom\Rightarrow\Delta\) & \(\bot,\Gamma\Rightarrow\Delta\) & none \\
Negation rules & \(\neg L,\neg R\) primitive & none & primitive \\
\(\to L\) & shared contexts & split contexts & shared contexts \\
Cut & split contexts & split contexts, plus Mix & shared contexts \\
Cut-free result & admissibility, semantic proof (this section) &
 elimination by local reductions & admissibility, semantic proof \\
\bottomrule
\end{tabular}
\end{center}
The comparison \cref{thm:p07-lk-presentation-equivalence} between the
canonical primitive-negation calculus of Part~7 and
\texttt{sheet:lk-defneg} is proved there for derivability, and
\cref{annex:p7.calculi.2} states that it establishes no preservation of
Cut-freeness, height or size.  The Part~7 theorem therefore does not by itself give a
Cut-free \(\mathsf{LK}_{\mathsf{CPL}}\)-derivation, and the present theorem
does not depend on it.  Conversely, the present theorem says nothing about
quantifiers, Mix, uncontexted identity or split-context \(\to L\).  The
subformula statement \cref{thm:p02-w02b-p02-009-subformula} is the
\(\mathsf{LK}_{\mathsf{CPL}}\) counterpart of
the subformula clause of \cref{thm:p07-classical-analytic}; here negation is primitive, so no
expansion of \(\neg\) precedes the subformula test.
\end{hbconvention}

\subsection{From natural deduction to sequents}

\begin{hbdefinition}{def:p02-w02b-p02-009-open-multiset}{Open-assumption multiset}
For an \(\mathsf{ND}_{\mathsf{CPL}}\)-derivation \(D\), the
\emph{open-assumption multiset} \(\operatorname{om}(D)\) is the finite
multiset whose value at \(\chi\) is the number of open assumption
occurrences of \(\chi\) in \(D\).  Its support is the set of formulas having
an open occurrence in \(D\).  If \(r\) is the root rule of \(D\) and
\(\operatorname{Dis}(r)\) is the set of occurrences it discharges, then the
open occurrences of \(D\) are those of its immediate subderivations, minus
\(\operatorname{Dis}(r)\); the occurrence sets of distinct immediate
subderivations are disjoint, so their multisets add.
\end{hbdefinition}

\begin{hbmethod}{def:p02-w02b-p02-009-nd-to-lk}{Translating a natural-deduction derivation into a sequent derivation}
\emph{Input:} an \(\mathsf{ND}_{\mathsf{CPL}}\)-derivation \(D\) with
conclusion \(\varphi\).  \emph{Output:} an
\(\mathsf{LK}_{\mathsf{CPL}}\)-derivation \(D^{\mathsf{LK}}\) of
\(\operatorname{om}(D)\Rightarrow\varphi\).  Proceed by recursion on \(D\).
Write \(D_1,D_2,D_3\) for the immediate subderivations in left-to-right
order and \(\Gamma_m\defeq\operatorname{om}(D_m)\).  A label
\((W_L)^{\Pi}\) denotes \(|\Pi|\) successive \(W_L\) steps adding the members
of \(\Pi\) with their multiplicities, \((C_L)^{n}\) denotes \(n\) successive
\(C_L\) steps on the displayed formula, and either denotes no step when its
count is zero.  Each translation tree is written as numbered lines; a line
names the rule that concludes it and the lines of its premises, and a line
marked ``from \(D_m\)'' is the root of the translation of \(D_m\).
\begin{enumerate}
  \item \emph{Assumption} \(\varphi^{\ell}\) (open): the identity sequent
    \(\varphi\Rightarrow\varphi\).  It discharges nothing;
    \(\operatorname{om}(D)=\varphi\).
  \item \(\land I\):
    \[
    \begin{array}{rll}
    1.&\Gamma_1\Rightarrow\varphi&\text{from }D_1\\
    2.&\Gamma_2\Rightarrow\psi&\text{from }D_2\\
    3.&\Gamma_1,\Gamma_2\Rightarrow\varphi&(W_L)^{\Gamma_2},\ 1\\
    4.&\Gamma_1,\Gamma_2\Rightarrow\psi&(W_L)^{\Gamma_1},\ 2\\
    5.&\Gamma_1,\Gamma_2\Rightarrow\varphi\land\psi&\land R,\ 3,4
    \end{array}
    \]
  \item \(\land E_1\) (for \(\land E_2\) replace line 2 by the identity
    sequent \(\varphi,\psi\Rightarrow\psi\) and conclude \(\psi\)):
    \[
    \begin{array}{rll}
    1.&\Gamma_1\Rightarrow\varphi\land\psi&\text{from }D_1\\
    2.&\varphi,\psi\Rightarrow\varphi&\text{identity sequent}\\
    3.&\varphi\land\psi\Rightarrow\varphi&\land L,\ 2\\
    4.&\Gamma_1\Rightarrow\varphi&\mathrm{Cut},\ 1,3
    \end{array}
    \]
  \item \(\lor I_1\) (for \(\lor I_2\) start from \(\Gamma_1\Rightarrow\psi\)
    and let \(W_R\) add \(\varphi\)):
    \[
    \begin{array}{rll}
    1.&\Gamma_1\Rightarrow\varphi&\text{from }D_1\\
    2.&\Gamma_1\Rightarrow\varphi,\psi&W_R,\ 1\\
    3.&\Gamma_1\Rightarrow\varphi\lor\psi&\lor R,\ 2
    \end{array}
    \]
  \item \(\to E\) (premises \(\varphi\) from \(D_1\), \(\varphi\to\psi\) from
    \(D_2\)):
    \[
    \begin{array}{rll}
    1.&\Gamma_1\Rightarrow\varphi&\text{from }D_1\\
    2.&\Gamma_2\Rightarrow\varphi\to\psi&\text{from }D_2\\
    3.&\Gamma_1\Rightarrow\psi,\varphi&W_R,\ 1\\
    4.&\Gamma_1,\psi\Rightarrow\psi&\text{identity sequent}\\
    5.&\Gamma_1,\varphi\to\psi\Rightarrow\psi&\to L,\ 3,4\\
    6.&\Gamma_2,\Gamma_1\Rightarrow\psi&\mathrm{Cut},\ 2,5
    \end{array}
    \]
  \item \(\neg E\) (premises \(\varphi\) from \(D_1\), \(\neg\varphi\) from
    \(D_2\)):
    \[
    \begin{array}{rll}
    1.&\Gamma_1\Rightarrow\varphi&\text{from }D_1\\
    2.&\Gamma_2\Rightarrow\neg\varphi&\text{from }D_2\\
    3.&\Gamma_1,\neg\varphi\Rightarrow{}&\neg L,\ 1\\
    4.&\Gamma_1,\neg\varphi\Rightarrow\bottom&W_R,\ 3\\
    5.&\Gamma_2,\Gamma_1\Rightarrow\bottom&\mathrm{Cut},\ 2,4
    \end{array}
    \]
  \item \(\bottom E\):
    \[
    \begin{array}{rll}
    1.&\Gamma_1\Rightarrow\bottom&\text{from }D_1\\
    2.&\bottom\Rightarrow\varphi&\bottom\text{-sequent}\\
    3.&\Gamma_1\Rightarrow\varphi&\mathrm{Cut},\ 1,2
    \end{array}
    \]
  \item \(\to I\) with label \(i\): let \(k\geq1\) be the number of open
    occurrences of \(\varphi\) marked \(i\) in \(D_1\); these are exactly the
    discharged ones, and \(\Gamma_1=\Gamma',\varphi^{k}\) with
    \(\Gamma'=\operatorname{om}(D)\):
    \[
    \begin{array}{rll}
    1.&\Gamma',\varphi^{k}\Rightarrow\psi&\text{from }D_1\\
    2.&\Gamma',\varphi\Rightarrow\psi&(C_L)^{k-1},\ 1\\
    3.&\Gamma'\Rightarrow\varphi\to\psi&\to R,\ 2
    \end{array}
    \]
  \item \(\neg I\) with label \(i\), \(k\) and \(\Gamma'\) as in item 8:
    \[
    \begin{array}{rll}
    1.&\Gamma',\varphi^{k}\Rightarrow\bottom&\text{from }D_1\\
    2.&\Gamma',\varphi\Rightarrow\bottom&(C_L)^{k-1},\ 1\\
    3.&\bottom\Rightarrow{}&\bottom\text{-sequent}\\
    4.&\Gamma',\varphi\Rightarrow{}&\mathrm{Cut},\ 2,3\\
    5.&\Gamma'\Rightarrow\neg\varphi&\neg R,\ 4
    \end{array}
    \]
  \item \(\mathsf{RAA}\) with label \(i\): let \(k\geq1\) be the number of
    open occurrences of \(\neg\varphi\) marked \(i\) in \(D_1\), and
    \(\Gamma_1=\Gamma',(\neg\varphi)^{k}\):
    \[
    \begin{array}{rll}
    1.&\Gamma',(\neg\varphi)^{k}\Rightarrow\bottom&\text{from }D_1\\
    2.&\Gamma',\neg\varphi\Rightarrow\bottom&(C_L)^{k-1},\ 1\\
    3.&\bottom\Rightarrow{}&\bottom\text{-sequent}\\
    4.&\Gamma',\neg\varphi\Rightarrow{}&\mathrm{Cut},\ 2,3\\
    5.&\varphi\Rightarrow\varphi&\text{identity sequent}\\
    6.&{}\Rightarrow\varphi,\neg\varphi&\neg R,\ 5\\
    7.&\Gamma'\Rightarrow\varphi&\mathrm{Cut},\ 6,4
    \end{array}
    \]
  \item \(\lor E\) with labels \(i,j\) and conclusion \(\chi\): let
    \(k\geq1\) be the number of open occurrences of \(\varphi\) marked
    \(i\) in \(D_2\) and \(m\geq1\) that of \(\psi\) marked \(j\) in \(D_3\);
    write \(\Gamma_2=\Gamma_2',\varphi^{k}\) and
    \(\Gamma_3=\Gamma_3',\psi^{m}\):
    \[
    \begin{array}{rll}
    1.&\Gamma_1\Rightarrow\varphi\lor\psi&\text{from }D_1\\
    2.&\Gamma_2',\varphi^{k}\Rightarrow\chi&\text{from }D_2\\
    3.&\Gamma_3',\psi^{m}\Rightarrow\chi&\text{from }D_3\\
    4.&\Gamma_2',\varphi\Rightarrow\chi&(C_L)^{k-1},\ 2\\
    5.&\Gamma_2',\Gamma_3',\varphi\Rightarrow\chi&(W_L)^{\Gamma_3'},\ 4\\
    6.&\Gamma_3',\psi\Rightarrow\chi&(C_L)^{m-1},\ 3\\
    7.&\Gamma_2',\Gamma_3',\psi\Rightarrow\chi&(W_L)^{\Gamma_2'},\ 6\\
    8.&\Gamma_2',\Gamma_3',\varphi\lor\psi\Rightarrow\chi&\lor L,\ 5,7\\
    9.&\Gamma_1,\Gamma_2',\Gamma_3'\Rightarrow\chi&\mathrm{Cut},\ 1,8
    \end{array}
    \]
\end{enumerate}
\emph{Side conditions:} the counts \(k,m\) are at least \(1\) because every
label of a discharging application discharges at least one occurrence;
occurrences marked \(i\) of a formula other than the bracketed one, and
occurrences marked \(i\) above another premise of \(\lor E\), are not
discharged and remain in \(\Gamma'\), \(\Gamma_1\), \(\Gamma_2'\) or
\(\Gamma_3'\).
\end{hbmethod}

\begin{hbtheorem}{thm:p02-w02b-p02-009-nd-to-lk}{proved=annex:p02-w02b-p02-009-nd-to-lk}{Natural deduction into the sequent calculus}
For every \(\mathsf{ND}_{\mathsf{CPL}}\)-derivation \(D\) with conclusion
\(\varphi\), the tree \(D^{\mathsf{LK}}\) of
\cref{def:p02-w02b-p02-009-nd-to-lk} is an
\(\mathsf{LK}_{\mathsf{CPL}}\)-derivation of
\(\operatorname{om}(D)\Rightarrow\varphi\).  In it each rule application
discharging \(k\) occurrences produces exactly \(k-1\) applications of
\(C_L\), and each two- or three-premise rule produces exactly
\(|\Gamma_2|+|\Gamma_1|\) (for \(\land I\)) or
\(|\Gamma_3'|+|\Gamma_2'|\) (for \(\lor E\)) applications of \(W_L\).
\end{hbtheorem}

\subsection{From sequents to natural deduction}

\begin{hbdefinition}{def:p02-w02b-p02-009-representation}{Label scheme and representing derivations}
The set \(\mathsf{Fm}_{\mathsf{CPL}}\) consists of finite strings over a
countable alphabet and is therefore countable.  Fix maps
\(\ell_{\mathrm L},\ell_{\mathrm R},\ell_{\mathrm A}:
\mathsf{Fm}_{\mathsf{CPL}}\to\N\) that are jointly injective: the values
\(\ell_{\mathrm X}(\chi)\) for \(\mathrm X\in\{\mathrm L,\mathrm R,\mathrm A\}\)
and \(\chi\in\mathsf{Fm}_{\mathsf{CPL}}\) are pairwise distinct.  An
\(\mathsf{ND}_{\mathsf{CPL}}\)-derivation \(D\) \emph{represents} the
sequent \(\Gamma\Rightarrow\Delta\) iff its conclusion is \(\bottom\) and each
of its open assumption occurrences is either an occurrence of some
\(\gamma\in\operatorname{supp}(\Gamma)\) marked \(\ell_{\mathrm L}(\gamma)\)
(a \emph{left} occurrence) or an occurrence of \(\neg\delta\) for some
\(\delta\in\operatorname{supp}(\Delta)\) marked \(\ell_{\mathrm R}(\delta)\)
(a \emph{right} occurrence).  Not every member of the supports need occur.
\end{hbdefinition}

\begin{hbmethod}{def:p02-w02b-p02-009-lk-to-nd}{Translating a sequent derivation into a natural-deduction derivation}
\emph{Tree notation.}  For a rule \(R\) of \(\mathsf{ND}_{\mathsf{CPL}}\),
\(R(T_1,\ldots,T_n)\) is the tree whose root is an application of \(R\)
with immediate subtrees \(T_1,\ldots,T_n\), in the left-to-right order of the
premises of \(R\) as displayed with the rules; \(R[\ell](\ldots)\) and
\(R[\ell,\ell'](\ldots)\) record the labels carried by a discharging
application; a marked leaf is \(\chi^{\ell}\).  In \(\lor I_1\) and
\(\lor I_2\) below the conclusion is the displayed disjunction.

\emph{Helper operations} on a derivation \(D\) with conclusion \(\bottom\):
\begin{itemize}
  \item \(\operatorname{ext}_\chi(D)\), concluding \(\chi\): if \(D\) has an
    open occurrence of \(\neg\chi\) marked \(\ell_{\mathrm R}(\chi)\), it is
    \(\mathsf{RAA}[\ell_{\mathrm R}(\chi)](D)\), discharging exactly those
    occurrences; otherwise it is \(\bottom E(D)\), discharging nothing.
  \item \(\operatorname{neg}_\chi(D)\), concluding \(\neg\chi\): if \(D\)
    has an open occurrence of \(\chi\) marked \(\ell_{\mathrm L}(\chi)\), it
    is \(\neg I[\ell_{\mathrm L}(\chi)](D)\), discharging exactly those
    occurrences; otherwise it is \(\bottom E(D)\).
  \item \(\operatorname{pin}_\chi(D)\defeq
    \neg E\bigl(\chi^{\ell_{\mathrm A}(\chi)},\operatorname{neg}_\chi(D)\bigr)\),
    concluding \(\bottom\), with one new open occurrence of \(\chi\) marked
    \(\ell_{\mathrm A}(\chi)\).
\end{itemize}
\emph{Input:} an \(\mathsf{LK}_{\mathsf{CPL}}\)-derivation of
\(\Gamma\Rightarrow\Delta\).  \emph{Output:} a derivation \(D^{\mathsf{ND}}\)
representing it.  Proceed by recursion on the sequent derivation; \(D_1,D_2\)
are the outputs for the premises.
\begin{enumerate}
  \item Identity \(\Gamma,\varphi\Rightarrow\varphi,\Delta\):
    \[\neg E\bigl(\varphi^{\ell_{\mathrm L}(\varphi)},
    (\neg\varphi)^{\ell_{\mathrm R}(\varphi)}\bigr)\]
  \item \(\bottom\)-sequent \(\Gamma,\bottom\Rightarrow\Delta\): the one-node
    tree \(\bottom^{\ell_{\mathrm L}(\bottom)}\).
  \item \(W_L,W_R,C_L,C_R\): \(D_1\) itself.
  \item \(\neg L\) (premise \(\Gamma\Rightarrow\Delta,\varphi\)):
    \[\neg E\bigl(\operatorname{ext}_\varphi(D_1),
    (\neg\varphi)^{\ell_{\mathrm L}(\neg\varphi)}\bigr)\]
  \item \(\neg R\) (premise \(\Gamma,\varphi\Rightarrow\Delta\)):
    \[\neg E\bigl(\operatorname{neg}_\varphi(D_1),
    (\neg\neg\varphi)^{\ell_{\mathrm R}(\neg\varphi)}\bigr)\]
  \item \(\land L\) (premise \(\Gamma,\varphi,\psi\Rightarrow\Delta\)): first
    \[D'\defeq\neg E\bigl(\land E_2((\varphi\land\psi)^{\ell_{\mathrm L}(\varphi\land\psi)}),
    \operatorname{neg}_\psi(D_1)\bigr)\]
    and then
    \[\neg E\bigl(\land E_1((\varphi\land\psi)^{\ell_{\mathrm L}(\varphi\land\psi)}),
    \operatorname{neg}_\varphi(D')\bigr)\]
  \item \(\land R\) (premises \(\Gamma\Rightarrow\Delta,\varphi\) and
    \(\Gamma\Rightarrow\Delta,\psi\)):
    \[\neg E\bigl(\land I(\operatorname{ext}_\varphi(D_1),\operatorname{ext}_\psi(D_2)),
    (\neg(\varphi\land\psi))^{\ell_{\mathrm R}(\varphi\land\psi)}\bigr)\]
  \item \(\lor L\) (premises \(\Gamma,\varphi\Rightarrow\Delta\) and
    \(\Gamma,\psi\Rightarrow\Delta\)):
    \[\lor E[\ell_{\mathrm A}(\varphi),\ell_{\mathrm A}(\psi)]\bigl(
    (\varphi\lor\psi)^{\ell_{\mathrm L}(\varphi\lor\psi)},
    \operatorname{pin}_\varphi(D_1),\operatorname{pin}_\psi(D_2)\bigr)\]
    Its conclusion is \(\bottom\).
  \item \(\lor R\) (premise \(\Gamma\Rightarrow\Delta,\varphi,\psi\)): first
    \[D'\defeq\neg E\bigl(\lor I_2(\operatorname{ext}_\psi(D_1)),
    (\neg(\varphi\lor\psi))^{\ell_{\mathrm R}(\varphi\lor\psi)}\bigr)\]
    and then
    \[\neg E\bigl(\lor I_1(\operatorname{ext}_\varphi(D')),
    (\neg(\varphi\lor\psi))^{\ell_{\mathrm R}(\varphi\lor\psi)}\bigr)\]
  \item \(\to L\) (premises \(\Gamma\Rightarrow\Delta,\varphi\) and
    \(\Gamma,\psi\Rightarrow\Delta\)):
    \[\neg E\bigl(\to E(\operatorname{ext}_\varphi(D_1),
    (\varphi\to\psi)^{\ell_{\mathrm L}(\varphi\to\psi)}),
    \operatorname{neg}_\psi(D_2)\bigr)\]
  \item \(\to R\) (premise \(\Gamma,\varphi\Rightarrow\Delta,\psi\)):
    \[\neg E\bigl(\to I[\ell_{\mathrm A}(\varphi)](
    \operatorname{ext}_\psi(\operatorname{pin}_\varphi(D_1))),
    (\neg(\varphi\to\psi))^{\ell_{\mathrm R}(\varphi\to\psi)}\bigr)\]
  \item Cut (premises \(\Gamma\Rightarrow\Delta,\varphi\) and
    \(\Sigma,\varphi\Rightarrow\Pi\)):
    \[\neg E\bigl(\operatorname{ext}_\varphi(D_1),
    \operatorname{neg}_\varphi(D_2)\bigr)\]
\end{enumerate}
\emph{Side conditions:} every discharging application above is nonvacuous:
\(\mathsf{RAA}\) and \(\neg I\) are used only when an occurrence is present,
and the auxiliary discharges in items 8 and 11 each discharge exactly the
single occurrence created by \(\operatorname{pin}\).  No subderivation is
grafted onto a leaf, so no label is renamed.
\end{hbmethod}

\begin{hbtheorem}{thm:p02-w02b-p02-009-lk-to-nd}{proved=annex:p02-w02b-p02-009-lk-to-nd}{Sequent derivations into natural deduction}
For every \(\mathsf{LK}_{\mathsf{CPL}}\)-derivation, with or without Cut, of
\(\Gamma\Rightarrow\Delta\), the tree \(D^{\mathsf{ND}}\) of
\cref{def:p02-w02b-p02-009-lk-to-nd} is an
\(\mathsf{ND}_{\mathsf{CPL}}\)-derivation representing
\(\Gamma\Rightarrow\Delta\).  At each step the discharged occurrences are
exactly: for \(\operatorname{ext}_\chi\), the open right occurrences of
\(\neg\chi\) marked \(\ell_{\mathrm R}(\chi)\); for
\(\operatorname{neg}_\chi\), the open left occurrences of \(\chi\) marked
\(\ell_{\mathrm L}(\chi)\); for the auxiliary labels, the one occurrence
created by \(\operatorname{pin}\).  Weakening and contraction leave the tree
unchanged: representation depends only on the supports of the two sides,
which weakening enlarges and contraction preserves.
\end{hbtheorem}

\begin{hbcorollary}{thm:p02-w02b-p02-009-equivalence}{proved=annex:p02-w02b-p02-009-equivalence}{Derivability correspondence by explicit translation}
Let \(\Gamma\) be a finite set of formulas, read as a multiset with
multiplicity one, and \(\Delta\) a finite multiset.
\begin{enumerate}
  \item If \(\Gamma\derives_{\mathsf{ND}_{\mathsf{CPL}}}\varphi\), then
    \(\Gamma\Rightarrow\varphi\) is derivable in \(\mathsf{LK}_{\mathsf{CPL}}\),
    and it has a purely logical derivation.
  \item If \(\Gamma\Rightarrow\Delta\) is derivable in
    \(\mathsf{LK}_{\mathsf{CPL}}\), then
    \(\Gamma\cup\{\neg\delta:\delta\in\operatorname{supp}(\Delta)\}
    \derives_{\mathsf{ND}_{\mathsf{CPL}}}\bottom\); if \(\Delta=\varphi\), then
    \(\Gamma\derives_{\mathsf{ND}_{\mathsf{CPL}}}\varphi\).
\end{enumerate}
Apart from the words ``purely logical'' in item 1, which use
\cref{thm:p02-w02b-p02-009-cut-admissibility}, the proof is syntactic: it
uses neither soundness nor completeness.
\end{hbcorollary}

\begin{hbexample}{def:p02-w02b-p02-009-example-nd-lk}{From natural deduction to sequents}
Take the derivation of \(P\lor\neg P\) in the excluded-middle example
before \cref{lem:nd-substitution}, with two occurrences of
\(\neg(P\lor\neg P)\) marked \(1\) and one occurrence of \(P\) marked \(2\).
Write \(\theta\) for \(P\lor\neg P\).  The translation produces, in order:
\begin{enumerate}
  \item \(P\Rightarrow P\) (assumption \(P^{2}\)); \(W_R\) and \(\lor R\)
    give \(P\Rightarrow\theta\) (item 4 of the method).
  \item The rule \(\neg E\) with the assumption \((\neg\theta)^{1}\),
    translated as \(\neg\theta\Rightarrow\neg\theta\): \(\neg L\) and
    \(W_R\) give \(P,\neg\theta\Rightarrow\bottom\), and Cut against
    \(\neg\theta\Rightarrow\neg\theta\) gives
    \(\neg\theta,P\Rightarrow\bottom\).  The antecedent records the two
    open occurrences.
  \item \(\neg I\) with label \(2\) discharges one occurrence (\(k=1\), no
    contraction): Cut against \(\bottom\Rightarrow{}\) and \(\neg R\) give
    \(\neg\theta\Rightarrow\neg P\).  Then \(W_R\) and \(\lor R\) give
    \(\neg\theta\Rightarrow\theta\).
  \item The rule \(\neg E\) with the second assumption \((\neg\theta)^{1}\):
    as in step 2, \(\neg\theta,\neg\theta\Rightarrow\bottom\).  Its
    antecedent has multiplicity \(2\), matching the two open occurrences
    marked \(1\).
  \item \(\mathsf{RAA}\) with label \(1\) discharges both (\(k=2\)): one
    \(C_L\) gives \(\neg\theta\Rightarrow\bottom\); Cut against
    \(\bottom\Rightarrow{}\) gives \(\neg\theta\Rightarrow{}\); Cut against
    \({}\Rightarrow\theta,\neg\theta\), obtained from
    \(\theta\Rightarrow\theta\) by \(\neg R\), gives \({}\Rightarrow\theta\).
\end{enumerate}
The result derives \({}\Rightarrow P\lor\neg P\) with five Cuts.  The
purely logical derivation of the same sequent found by
\cref{def:p02-w02b-p02-009-decomposition} is
\[
\begin{array}{rll}
1.&P\Rightarrow P&\text{identity sequent}\\
2.&{}\Rightarrow P,\neg P&\neg R,\ 1\\
3.&{}\Rightarrow P\lor\neg P&\lor R,\ 2
\end{array}
\]
\end{hbexample}

\begin{hbexample}{def:p02-w02b-p02-009-example-lk-nd}{From sequents to natural deduction}
Translate the derivation in the example after the rules of
\(\mathsf{LK}_{\mathsf{CPL}}\): \(P\Rightarrow P,Q\), then \(\lor R\), then
\(\to R\).  Choose the label maps so that \(\ell_{\mathrm L}(P)=1\), \(\ell_{\mathrm R}(P)=2\),
\(\ell_{\mathrm R}(P\lor Q)=3\), \(\ell_{\mathrm A}(P)=4\),
\(\ell_{\mathrm R}(P\to(P\lor Q))=5\).
\[
\begin{array}{rll}
1.&P^{1}&\text{assumption}\\
2.&\neg P^{2}&\text{assumption}\\
3.&\bottom&\neg E,1,2\quad\text{(identity \(P\Rightarrow P,Q\))}\\
4.&Q&\bottom E,3\quad\text{(\(\operatorname{ext}_Q\): no \(\neg Q\) marked \(\ell_{\mathrm R}(Q)\))}\\
5.&P\lor Q&\lor I_2,4\\
6.&\neg(P\lor Q)^{3}&\text{assumption}\\
7.&\bottom&\neg E,5,6\\
8.&P&\mathsf{RAA},2\text{--}7\quad\text{(discharges line 2)}\\
9.&P\lor Q&\lor I_1,8\\
10.&\neg(P\lor Q)^{3}&\text{assumption}\\
11.&\bottom&\neg E,9,10\quad\text{(represents \(P\Rightarrow P\lor Q\))}\\
12.&\neg P&\neg I,1\text{--}11\quad\text{(discharges line 1)}\\
13.&P^{4}&\text{assumption}\\
14.&\bottom&\neg E,13,12\\
15.&P\lor Q&\mathsf{RAA},3\text{--}14\quad\text{(discharges lines 6 and 10)}\\
16.&P\to(P\lor Q)&\to I,4\text{--}15\quad\text{(discharges line 13)}\\
17.&\neg(P\to(P\lor Q))^{5}&\text{assumption}\\
18.&\bottom&\neg E,16,17\quad\text{(represents \({}\Rightarrow P\to(P\lor Q)\))}\\
19.&P\to(P\lor Q)&\mathsf{RAA},5\text{--}18.
\end{array}
\]
Line 19 is \(\operatorname{ext}_{P\to(P\lor Q)}\) of line 18 and has no open
assumptions.
\end{hbexample}

\begin{hbwarning}{def:p02-w02b-p02-009-translation-limits}{What the translations preserve}
The two translations are not inverse to each other and preserve neither
size nor the absence of detours.  The translation into sequents uses Cut for
every elimination rule and for \(\neg I\) and \(\mathsf{RAA}\), so it does
not produce Cut-free derivations; Cut-freeness is recovered only through
\cref{thm:p02-w02b-p02-009-cut-admissibility}, which does not start from the
translated derivation.  The translation into natural deduction turns every
Cut, and every right rule, into an introduction followed by an elimination
or by \(\neg E\) against a negated assumption.  It forgets multiplicities:
contraction and weakening act as the identity on trees, and the number of
open occurrences of a formula in \(D^{\mathsf{ND}}\) is not determined by
its multiplicity in the sequent.  Natural-deduction derivability depends only on
the set of open formulas; sequent derivability depends on the multisets, but
only up to the structural rules.
\end{hbwarning}

\paragraph{Relationships and sources.}
The equivalence of the three calculi with the semantic consequence relation
is \cref{thm:p02-calculi-derivability-equivalence}; the translations above
give its natural-deduction and sequent legs by explicit constructions on
\(\mathsf{ND}_{\mathsf{CPL}}\) and \(\mathsf{LK}_{\mathsf{CPL}}\) themselves,
in place of the derivability-level construction of
\cref{thm:p07-lk-nd-equivalence} for the first-order calculi.  The sketched
completeness argument \cref{thm:prop-lk-completeness} uses the
natural-deduction-to-sequent direction; its translation subroutine is
\cref{thm:p02-w02b-p02-009-nd-to-lk}, and
\cref{thm:p02-w02b-p02-009-cutfree-completeness} gives completeness without
Cut by a separate argument.  The semantic route to Cut-free completeness and
Cut elimination, the inversion property and the subformula property for a
propositional sequent calculus are in Buss
\cite[Secs.~1.2.4--1.2.11, pp.~13--16]{Buss1998Intro}.  The syntactic
elimination for the defined-negation calculus is
\cref{thm:p07-classical-cut-elimination}.
